\documentclass[10pt,a4paper]{amsart}
\usepackage[margin=19mm]{geometry}
\usepackage{amsmath,amssymb,mathtools}
\usepackage[scr=rsfs,cal=euler]{mathalfa}
\usepackage{xfrac}
\usepackage[unicode,hidelinks,bookmarks=false]{hyperref}
\newcommand{\ie}{i.e.\ }
\newcommand{\cf}{cf.\ }
\newcommand{\resp}{respectively\ }
\newcommand{\Subsection}{\subsection}
\providecommand{\nobreakdash}{\nobreak\hbox{-}}
\setlength{\emergencystretch}{2em}
\multlinegap0pt
\usepackage{amssymb}
\usepackage[shortlabels]{enumitem}
\usepackage{mathtools}
\usepackage{bm}
\newtheorem{thm}{Theorem}
\title{Some Generalizations of Totient Function with Elementary Symmetric Sums}
\author{Udvas Acharjee, N. Uday Kiran}
\date{Source: arXiv:2511.19502v2 (CC BY 4.0)}
\begin{document}
\maketitle

\noindent\emph{Source and licence.} Some Generalizations of Totient Function with Elementary Symmetric Sums. Version arXiv:2511.19502v2 (CC BY 4.0), distributed by arXiv under the Creative Commons Attribution 4.0 International licence, http://creativecommons.org/licenses/by/4.0/, as recorded in the arXiv licence metadata for this exact version and shown on the abstract page https://arxiv.org/abs/2511.19502v2. Full licence notice: \texttt{licenses/arxiv-2511.19502-CC-BY-4.0.txt}.\par\smallskip

\noindent\textit{Editorial scope.} Counted unit: the article's thm tot\_func2(1,2 (source0.tex lines 391--407). The article's complete original proof of it is delivered with it (source0.tex lines 408--418, one \texttt{proof} environment of 11 lines). No mathematics is written, repaired or re-derived: the statement and the proof are the article's own lines, in the article's own notation, and no retained line is substituted or re-flowed.  Retained uncounted as the source-local context the proof needs: lines 351--355.  Nothing was omitted from any retained span, so the package carries no omission record.  No retained line contains a citation, so the wrapper carries no bibliography and no bibliography entry.  No retained line contains a figure environment, an image-inclusion command or drawing code, so no image is delivered and the package declares no asset dependency.  Editorial formatting only, in addition: an AMS \texttt{amsart} preamble that restates the article's own declarations and macros used by the retained lines, namely the environments \texttt{thm} on separate counters, together with the standard packages those lines need.  The retained environments print in this document as Theorem 1 in order of appearance, whereas the article prints the same items with the numbering of its own full document.  The complete native source of the article as distributed in the arXiv e-print for this exact version is bundled unchanged as \texttt{sources/189722/source0.tex}.
\begin{thm}[Relation between totient functions]\label{thm: relation_tot}
\[
\phi_F(p^k)=\sum_{\substack{J \subseteq F \\ J \ne \varnothing}}(-1)^{|J|+1}\varphi_J(p^k), \quad \varphi_F(p^k)=\sum_{\substack{J \subseteq F \\ J \ne \varnothing}}(-1)^{|J|+1}\phi_J(p^k).
\]
\end{thm}

\begin{thm}\label{thm: tot_func2{1,2}}
    For $k\geq 2$, the following identity holds for odd $n$:\[
    \phi_{\{1,2\}}(n)=n^k\prod_{p\mid n}\left(1-\frac{1}{p}-\frac{p-1}{p^2}+(p-1)\frac{h_k(p)}{p^k}\right)
    \]
    where\[
    h_k(p)=\begin{cases}
        p^{(k-3)/2}\eta((-1)^{(k-1)/2}k)-p^{(k-1)/2}\eta((-1)^{(k)}(1-\gcd(k-1, n))), &\text{ if }k\text{ is odd}\\
        p^{(k-2)/2}\eta((-1)^{k/2}(1-\gcd(k,p)))-
         p^{(k-2)/2}\eta((-1)^{k/2+1}(k-1)), &\text{ if }k\text{ is even}.
    \end{cases}
    \]
    Here $\eta(\cdot)$ is the quadratic character$\text{ mod } p$.
     Also,
    \[
    \phi_{\{1,2\}}(2^l)=2^{lk}\left(\frac{1}{4}-\frac{1}{2\left(\sqrt{2}\right)^k}\sin\frac{k\pi}{4}\right)
    \]
\end{thm}

\begin{proof}
    This theorem is obtained using theorem (\ref{thm: relation_tot}) as:
    \[
    \phi_{\{1,2\}}(n)=\varphi_{\{1\}}(n)+\varphi_{\{2\}}(n)-\varphi_{\{1,2\}}(n).
    \]
    This gives us:
    \[
    \phi_{\{1,2\}}(n)=n^k\prod_{p\mid n}\left(1-\frac{1}{p}+\frac{N_k(\{1,2\},p)-N_k(\{2\},p)}{p^k}\right)
    \]
    We get the result upon substituting the respective values of $N_k(\{2\},p)$ and $N_k(\{1,2\},p)$.
\end{proof}

\end{document}
